\chapter{Algoritmo DRI-SL}
\label{ap:drisl}

\section{Objetivo e entradas}

Construir $L_0$ de tamanho alvo $I$ a partir do \textit{pool} não rotulado $U$, sem
acesso a rótulos. Entradas: textos de $U$; codificador de sentenças $\phi$
(SBERT multilíngue \citep{Reimers2019SBERT} na configuração da tese); número
de grupos $N_c$ (padrão $\max(2, \lfloor\sqrt{I}\rfloor)$); semente.

\section{Etapas}

\begin{enumerate}
  \item \textbf{Densidade semântica}: projetar $U$ por $\phi$; agrupar por
  $k$-médias em $N_c$ grupos; alocar a cota $q_c$ de cada grupo
  proporcionalmente ao seu tamanho (cota mínima de 1 para grupo não vazio;
  ajuste do resto pelas maiores frações).
  \item \textbf{Variedade lexical intragrupo}: para cada grupo $c$, computar o
  perfil TF--IDF médio; selecionar iterativamente, entre os membros ainda não
  escolhidos, o texto que maximiza o \emph{escore de novidade} (a soma dos
  pesos, no perfil do grupo, dos termos que o texto introduz e que ainda não
  estão cobertos pelos já selecionados); atualizar o conjunto de termos
  cobertos; repetir até preencher $q_c$.
\end{enumerate}

A seleção é determinística dada a semente (inicialização do $k$-médias). O
custo é dominado pela codificação ($O(|U|)$ passagens do codificador) e pelo
$k$-médias; a etapa lexical é $O(q_c \cdot |c|)$ por grupo, com operações
esparsas.

\section{Intuição}

A etapa 1 garante \emph{representatividade}: como o número de grupos é
$N_c = \max(2, \lfloor\sqrt{I}\rfloor)$ e cada grupo não vazio tem cota mínima
de um, nenhum agrupamento fica de fora do $L_0$. A etapa 2 \emph{reduz} a
redundância: dentro da região, evita quase-duplicatas (frequentes no domínio:
7,7\% de duplicatas exatas) e força a cobertura de vocabulário, o recurso que,
para um classificador por protótipo como o PVBin, constrói as fronteiras. A combinação explica o
resultado do Capítulo~\ref{ch:resultados-l0}: cobertura semântica dá o
esqueleto de classes; cobertura lexical dá o poder discriminativo por classe.

Implementação de referência:
\texttt{src/activelearning/adapters/strategies/drisl.py} (com codificador
alternativo TF--IDF+SVD para ambientes leves), testes em
\texttt{tests/unit/test\_drisl.py}.
